\section{Proofs for branching comparisons}
\label{app:branching}

\subsection{Small certificates and the eleven-node example}
\label{app:branching-small-certificates}

\begin{proof}[The five-node assertion in \cref{branching:ratio-range}]
Let $s$ be the breakpoint of a two-piece valid partition of $[L,U]$.
If the root split is $s$, the tree has three nodes. By reflection suppose
it is $y_1>s$. Its right child is valid. If its left child is invalid,
let $y_2$ be its selected minimizer. The crossing lemma in the proof of
\cref{branching:four-competitive} excludes $y_2>s$: it would force the
root to cross its own right endpoint. If $y_2=s$, both new children
are valid. Otherwise $y_2<s$, and $[L,y_2]$ is valid.

Suppose $[y_2,y_1]$ is invalid, witnessed by a point $z$ with
$M_\epsilon(z)<\alpha(z-y_2)(y_1-z)$. Optimality of $y_2$ in
$[L,y_1]$ gives
\[
 M_\epsilon(y_2)<\alpha\bigl[(z-y_2)(y_1-z)
 +(y_2-L)(y_1-y_2)-(z-L)(y_1-z)\bigr]
 =\alpha(y_2-L)(z-y_2).
\]
Validity of $[L,s]$ at $y_2$ therefore implies $z>s$.
Optimality of $y_1$ in the root gives
\[
 M_\epsilon(y_1)<\alpha\bigl[(z-y_2)(y_1-z)
 +(y_1-L)(U-y_1)-(z-L)(U-z)\bigr]
 =\alpha(y_1-z)[(U-y_1)-(y_2-L)].
\]
Validity of $[s,U]$ gives
$M_\epsilon(y_1)\ge\alpha(y_1-s)(U-y_1)>0$.
If the bracket on the right is nonpositive there is an immediate
contradiction; otherwise dropping $y_2-L>0$ implies $z<s$.
This contradiction proves the middle interval valid. There are at most
two splits and five nodes. The information pair below attains five.
\end{proof}

Here is an explicit example for the lower endpoint $11/5$. Take
$\alpha=1$, $\epsilon=1/100$, and root $[0,1]$. Interpolate $H$ linearly
between the knots below, giving it values $H(x)=x^2+M(x)$ at those knots,
and set $f(x)=H(x)-x^2-\epsilon$:
\[
\begin{array}{c|ccccccccc}
 x&0&1/16&3/16&7/16&1/2&9/16&13/16&15/16&1\\ \hline
 M(x)&41/1600&41/1600&823/20000&1/100&1/100&1/100&
                         823/20000&41/1600&41/1600
\end{array}
\]
On each affine piece of $H$, $f$ is concave. Its minimum is attained
at a knot, so $f_*=0$. The slopes of $H$ are nondecreasing, and the
shifted node function is $H(x)-(l+u)x+lu$, which is affine on each piece.
Direct exact substitution gives the following invalid nodes and unique
minimizers:
\[
\begin{array}{c|c|c}
 {[l,u]}&y&\min(H-(l+u)x+lu)\\ \hline
 {[0,1]}&1/2&-6/25\\
 {[0,1/2]}&3/16&-2791/160000\\
 {[3/16,1/2]}&7/16&-9/1600\\
 {[1/2,1]}&13/16&-2791/160000\\
 {[1/2,13/16]}&9/16&-9/1600
\end{array}
\]
The six resulting leaves are valid: the shifted function is affine
between knots, and its values at all contained knots and the leaf
endpoints are nonnegative. Thus every minimizer selection uses eleven
nodes. A valid three-piece certificate has breakpoints $2/5$ and $3/5$.
To verify its validity, again evaluate $H-(l+u)x+lu$ at its endpoints
and all contained knots. The side minima are $209/160000$, and the
middle minimum is zero at $1/2$.
No two-piece certificate is possible. Put $a=12209/30000$.
Testing its left interval at $3/16$ forces $s\le a$ when $s\ge3/16$;
when $s<3/16$ this conclusion holds automatically.
Testing its right interval at $13/16$ forces $s\ge1-a$ when
$s\le13/16$; otherwise this conclusion holds automatically.
Since $1-a>a$, the two requirements contradict each other.
Hence $N_{\rm rect}=3$ and $T/T_{\rm opt}=11/5$.

\subsection{Two indistinguishable root instances}
\label{app:branching-information}

\begin{proof}[Proof of \cref{branching:germ-obstruction}]
Set $\alpha=1$, $\epsilon=1/10000$, and
$C_{ab}(y)=(a+b)y-ab$. Define
\begin{align*}
 g_-(y)&=143/300+(4/5)(y-3/5),&
 g_+(y)&=143/300+(6/5)(y-3/5),\\
 e_0(y)&=\epsilon+y/6,&
 e_1(y)&=(5/3)y-2/3+\epsilon,\\
 H_A&=\max\{C_{0,1/3},C_{1/3,1},g_-,g_+,e_0,e_1\},\\
 H_B&=\max\{C_{0,3/5},C_{3/5,1},g_-,g_+,e_0,e_1\},\\
 f_i&=H_i-y^2-\epsilon\qquad(i\in\{A,B\}).
\end{align*}
Both $H_i$ are convex. On $[0,1/6]$, $e_0-y^2\ge\epsilon$;
on $[2/3,1]$, $e_1-y^2\ge\epsilon$. On the two intervening intervals
$[1/6,3/5]$ and $[3/5,2/3]$, use $g_--y^2$ and $g_+-y^2$
respectively. These are concave quadratics whose endpoint values exceed
$\epsilon$. Thus $f_i\ge0$. At zero and one, $H_i=e_0$ and $e_1$
respectively, so both objectives have minimum zero.

The root shifted function is $H_i-y$. The common lower envelope
$\max(g_-,g_+)-y$ has its unique minimum $-37/300$ at $3/5$,
and both $H_i$ equal this envelope there. Hence both root minimizers
are uniquely $3/5$, and both actual relaxation bounds are
$-37/300-\epsilon$. Strict comparisons of the displayed affine lines
show that both functions equal $e_0$ near zero, $e_1$ near one,
and $\max(g_-,g_+)$ near $3/5$. Their root germs and endpoint
neighborhoods therefore agree. All node underestimators are convex,
since $f_{[l,u]}=H_i-(l+u)y+lu-\epsilon$.

For each instance, domination of its two chord lines gives a valid
partition at its named breakpoint $c_A=1/3$ or $c_B=3/5$.
Those breakpoints are unique. At $y=1/200$, $H_i=C_{0,c_i}$
with strict domination over its other lines; at $y=19/20$,
$H_i=C_{c_i,1}$ with strict domination. If a two-piece partition
has breakpoint $s$, validity at the first test point forces $s\le c_i$
unless $s<1/200$, which already satisfies this inequality. The second
test forces $s\ge c_i$ unless $s>19/20$, which already satisfies it.
Thus $s=c_i$. The root is invalid. Root data force the same deterministic
cut on both instances, and at least one cut is wrong. That instance
needs at least five nodes, against an optimum of three. For a randomized
cut, the probabilities of cutting exactly $c_A$ and $c_B$ sum to at
most one. Some wrong-cut event has probability at least $1/2$, giving
expected count at least $3+2(1/2)=4$.

To obtain smooth nonanalytic functions, extend each maximum of lines
to the real line and convolve with the same symmetric nonnegative
$C^\infty$ kernel of integral one, positive inside its support
$[-\zeta,\zeta]$. Choose $\zeta>0$ small enough that the endpoint
neighborhoods and the two tight-chord test neighborhoods stay affine,
and that the common neighborhood of $3/5$ contains the convolution
window. Symmetry gives zero first moment, so convolution preserves
these affine pieces. Jensen's inequality for convex $H_i$ makes its
convolution no smaller than $H_i$. Nonnegativity, endpoint minima,
chord domination, and the distinct unique certificate cuts all persist.
Near $3/5$ the common function is a line of slope one plus
$(1/5)|y-3/5|$. Its symmetric convolution minus $y$ has derivative
strictly increasing through zero at $3/5$, hence a unique minimum
there. Convexity excludes any other minimum outside that neighborhood.
The functions still agree on the needed open neighborhoods, and all
root data still agree. The preceding information argument applies.
Each smoothed $H_i$ is affine on an endpoint neighborhood and nonaffine
near the smoothed V at $3/5$. If it were analytic, continuation would
force it to be affine throughout. Thus each smoothed $H_i$, and therefore
each smoothed $f_i=H_i-y^2-\epsilon$, is nonanalytic.
\end{proof}

\subsection{Repeated ambiguity across coordinates}
\label{app:branching-dimension}

\begin{proof}[Proof of \cref{branching:dimension-obstruction}]
Fix $\epsilon=1/100$, $c=1/(5n)$, and put
\[
 \delta=\frac{\epsilon}{4(n-1)},\quad
 \varphi=(n-1)c-\epsilon/2,\quad
 \varphi'=c-\frac{\epsilon}{2(n-1)}=\frac{\varphi}{n-1}.
\]
On $[0,1]$ define
\begin{align*}
 h_-(t)&=1/2-c+(1-\delta)(t-1/2),\\
 h_+(t)&=1/2-c+(1+\delta)(t-1/2),\\
 H_R(t)&=\max\{h_-,h_+,t/2+\varphi,3t/2-1/2+\varphi\},\\
 H_N(t)&=\max\{h_-,h_+,t/2-\varphi',3t/2-1/2-\varphi'\},\\
 R(t)&=H_R(t)-t^2,\qquad N(t)=H_N(t)-t^2,\\
 f_i(x)&=R(x_i)+\sum_{j\ne i}N(x_j),\qquad i=1,\ldots,n.
\end{align*}
The coordinate underestimators are maxima of affine functions after
subtracting $q_J$, and hence convex.

Both root coordinate relaxations have unique minimizer $1/2$ and value
$-c$. Indeed at $1/2$ the common lines dominate the floor lines:
the $R$ margin is $1/4+\epsilon/2-nc=11/200>0$, and the $N$ margin
is $1/4-\epsilon/(2(n-1))>0$. Around this point $H-t$ has left slope
$-\delta$ and right slope $\delta$. Convexity makes its minimum
unique. Thus $R=N$ near $1/2$.
Near zero and one the respective floor lines dominate. At zero their
strict margins above $h_-$ are
$nc-\epsilon/2-\delta/2>0$ for $R$ and
$3\epsilon/(8(n-1))>0$ for $N$. The symmetric comparison holds at
one. Consequently $R-N=\varphi+\varphi'$ near both endpoints.

The floor lines give $R\ge\varphi$ and $N\ge-\varphi'$ on both
halves: subtracting $t^2$ leaves the nonnegative products
$t(1/2-t)$ and $(t-1/2)(1-t)$. Equality holds at the endpoints.
Thus $\min R=\varphi$, $\min N=-\varphi'$, and $\min f_i=0$.
Furthermore the $R$ relaxation value on either half is exactly
$\varphi$, while on any interval containing zero or one the $N$
relaxation value is at most $-\varphi'$ by evaluation at that endpoint.
Cutting coordinate $i$ at $1/2$ therefore makes both children valid:
their shifted relaxation value is
$\varphi-(n-1)c+\epsilon=\epsilon/2$.
The root is invalid, since its shifted value is $-nc+\epsilon<0$.
Hence every $f_i$ has a three-node optimum.

Call a box a corner box if each coordinate interval contains a root
endpoint. Let $S$ be its set of coordinates already cut; the others
are free and still have interval $[0,1]$. For every possible rigid
label $i\notin S$, each used coordinate is an $N$ coordinate, while
each free coordinate has the same unique root minimizer and value
whether it is $R$ or $N$. Coordinate-wise selection therefore gives
the same full minimizer and bound for all possible labels. At the
minimizer the free coordinates lie at $1/2$, where $R=N$, so the
objective germs also agree. Near each corner, the difference between
two labelled objectives is
$(R-N)(x_i)-(R-N)(x_{i'})=0$, because both differences are the same
constant at root endpoints. The corner germs agree as well.
Every such corner box is invalid: if $d=|S|\le n-1$, its shifted
value is at most
\[
 -(n-d)c-d\varphi'+\epsilon
 =-nc+\epsilon+\frac{d\epsilon}{2(n-1)}
 \le-1/5+3\epsilon/2<0.
\]

Fix a deterministic local rule. Until the rigid label has been cut,
the common corner-box data force the same decision for every remaining
label. These decisions define a reference tree of corner boxes.
Cutting a free coordinate creates two corner children and removes
that label from the free set. Re-cutting a used coordinate has one
corner child, with the same free set, and one noncorner child.
If such a re-cut chain never ends while a free label remains, that
instance has an infinite tree and satisfies the conclusion. Otherwise
contract the re-cut chains. The contracted reference tree has at least
$2^d$ branching nodes after $d$ free coordinates have been cut.

Let $I_i$ count reference corner nodes at which label $i$ is still
free. Every such node is reached and is internal in the actual tree
for $f_i$, by induction and the invalidity just proved. Thus
$T_i\ge2I_i+1$. Each branching node at level $d$ contributes to
$n-d$ labels, so
\[
 \sum_iI_i\ge\sum_{d=0}^{n-1}2^d(n-d)=2^{n+1}-n-2.
\]
If the root cuts coordinate $j$, then $I_j=1$. Averaging over the
other $n-1$ labels gives the deterministic bound.
For a distribution of local rules, average this last sum inequality
over the random rule and then over the $n$ labels. At least one label
has expected count at least $2(2^{n+1}-n-2)/n+1$.
The reference construction uses node-locality: objective-dependent
information from another subtree could identify the rigid label and
invalidate the repeated indistinguishability argument.
\end{proof}

\subsection{Safety lower bounds and recentring}
\label{app:branching-safety}

For the kink objective on $[0,1]$, the derivative of the underestimator
on the left of $a$ is $-2\alpha+\alpha(2y-l-u)<0$, and on its right
is $2\alpha+\alpha(2y-l-u)>0$. These signs also show that a one-sided
interval has its minimizer at the endpoint nearest $a$ and a
nonnegative value. Thus the invalidity test and off-chain validity used
in \cref{branching:persistent-clamp,branching:recentring} are exact.

\begin{proof}[Proof of \cref{branching:recentring}]
For $a<\theta$, every one of the first $J$ cuts on its containing
chain lies to its right. Indeed after $k$ cuts the left endpoint is zero
and the upper endpoint is at least $\theta^k$. For $k<J$,
$\theta^{k+1}>a$, so every admissible cut exceeds $a$. Also
$(a-l)(u-a)>a^2(1-\theta)/\theta$ at those nodes. The tolerance
hypothesis makes each node invalid, proving $T\ge2J+1$. A chain
that hits $a$ needs one further split. This proof holds samplewise for
randomized rules.
Choosing
$a=\sqrt{2\theta\epsilon/(\alpha(1-\theta))}<\theta$
for sufficiently small $\epsilon$ makes the node count logarithmic in
$1/\epsilon$, while $T_{\rm opt}=3$.

The recentring rule is safe: its offset from the nearby endpoint is
between $\theta w$ and $2\theta w$, and
$2\theta\le1-\theta$. Let
$d=\min\{(a-l)/w,(u-a)/w\}$. If $d<\theta/2$, an ordinary clamp
leaves a containing child of width $\theta w$ and new distance
$d/\theta<1/2$. If $\theta/2\le d<\theta$, the offset $2dw$
makes the kink the midpoint of its containing child. If $d\ge\theta$,
the next cut hits the kink.
Put $m=\min\{j\ge0:d_0\ge\theta^{j+1}/2\}$.
After $m$ ordinary clamps the distance is at least $\theta/2$.
If it is at least $\theta$, then $m=J$ and one more cut hits the kink.
Otherwise $J=m+1$: $d_0<\theta^{m+1}$, but
$d_0\ge\theta^{m+1}/2\ge\theta^{m+2}$. One recentring cut and one
exact kink cut finish. In either case the zero-tolerance chain takes
exactly $J+1$ splits. The preceding safety argument, reflected if
necessary, makes this the minimum hitting count. For any positive
tolerance the chain can stop earlier, giving $T\le2J+3$; for every
fixed kink equality holds at all sufficiently small positive tolerances,
because its finitely many pre-hit gap products are positive.
\end{proof}

\subsection{Why splitting all coordinates can overrefine}
\label{app:branching-multi}

\begin{proposition}[A two-dimensional multisection separation]
\label{branching:multi-separation}
For $f(x,z)=x^2$ on $[0,1]^2$ with $\alpha=1$, multisection at all
interior exact coordinate minimizers has
$\Omega(\epsilon^{-1/2}\log(1/\epsilon))$ internal nodes, whereas
$N_{\rm tree}=O(\epsilon^{-1/2})$. Count every processed box, whether
a split has two or four children.
\end{proposition}

\begin{proof}
The coordinate relaxation values are
\[
 F_x([l,u])=\begin{cases}lu-(l+u)^2/8,&u>3l,\\l^2,&u\le3l,
 \end{cases}\qquad F_z([c,d])=-(d-c)^2/4.
\]
The selected $x$ minimizer is $(l+u)/4$ in the first case and $l$
otherwise; the $z$ minimizer is always the midpoint.

For an upper certificate put $h=2\sqrt\epsilon$ and use $x$ strips
$[0,h]$, $[h,2h]$, $[2h,4h]$, and so on, truncating at one.
The first strip has shifted $x$ value $\epsilon/2$, so $z$ intervals
of width at most $\sqrt{2\epsilon}$ suffice. Each later strip
$[l,u]$ has $u\le2l$, hence $F_x=l^2$, and $z$ intervals of width
at most $2l$ suffice. Their counts sum to at most
$\lceil1/\sqrt{2\epsilon}\rceil+1/(2\sqrt\epsilon)
+O(\log(1/\epsilon))=O(\epsilon^{-1/2})$. These cuts form a
guillotine certificate.

In the multisection tree every depth-$D$ box has a $z$ interval of
width $2^{-D}$. For each $x$ history all $2^D$ such intervals occur
together, since validity depends only on the $z$ width.
The left $x$ spine is $[0,4^{-k}]$. Its right sibling
$J_k=[l,4l]$, $l=4^{-(k+1)}$, appears at depth $k+1$.
Inside $J_k$, let $r_0=l$ and $r_{j+1}=(r_j+4l)/4$.
The right descendants $[r_j,4l]$ split at $r_{j+1}$, with
$r_j\uparrow4l/3$. Their left pieces
$L_j=[r_{j-1},r_j]$ appear at depth $k+1+j$ and are never split
in $x$ again. Their shifted $x$ values are strictly below
$\epsilon+16^{-k}/9$. At depth $D\le2k$, the $z$ value is at most
$-16^{-k}/4$. All these boxes are therefore invalid when
$\epsilon\le(5/36)16^{-k}$; the strict $x$ estimate handles equality.
Hereditary validity ensures that their ancestors are also invalid, so
the indicated boxes really occur.
With $K=\lfloor\log_{16}(5/(36\epsilon))\rfloor$, count only depth
$2k$ and $j=1,\ldots,k-1$. These disjoint histories give at least
\[
 \sum_{k=2}^K(k-1)4^k\ge(K-1)4^K
       =\Omega(\epsilon^{-1/2}\log(1/\epsilon))
\]
internal nodes as $\epsilon\downarrow0$. A tree with at least two
children per internal node has at least $2I+1$ total nodes. The binary
certificate above has $2N-1$ nodes, proving the unbounded ratio.
\end{proof}

\subsection{A phase bound and the sharp-coordinate induction}
\label{app:branching-sharp}

Fix a continuous $m\ge0$ on an interval $A$, with a fixed exact
minimizer selection for every subinterval. Define
$F(J)=\min_J(m-q_J)$ and $\omega(J)=q_J(y_J)$.
The one-dimensional split-count argument applies unchanged to a
continuous $m\ge0$ at any budget $b>0$: it needs only a positive shifted
objective, and does not require a zero of $m$ on the current interval.
Let $\tau>0$, let $\kappa\ge1$ be an integer, take $\beta\in\R$, and put
$b=\beta+\kappa\tau>0$. Define a pruned minimizer tree $S$ that
splits $J$ precisely when $F(J)<-\beta$ and $\omega(J)\ge\tau$.

\begin{lemma}[Phase count]
\label{branching:phase-lemma}
Let $N=N_A(b)$, and put $c_\kappa=\kappa[2\kappa(\kappa+1)+1]$.
The number of internal nodes of $S$ is at most
$4N-5+c_\kappa(4N-4)$ when $N\ge2$, and at most $c_\kappa$
when $N=1$.
\end{lemma}

\begin{proof}
First take an interval $\ell=[p_0,q_0]$ with $F(\ell)\ge-b$.
For a split node $J=[p,q]\subseteq\ell$ with selected point $y$,
write $L=y-p$, $R=q-y$, $g_L=p-p_0$, $g_R=q_0-q$.
Validity of $\ell$ at budget $b$ and the split conditions imply
\begin{equation}
 Lg_R+Rg_L+g_Lg_R<\kappa\tau\le\kappa LR.
 \label{branching:phase-gap}
\end{equation}
Thus $g_R<\kappa R$, $g_L<\kappa L$, and
$|J|>|\ell|/(\kappa+1)$.
On a nested chain of split nodes contained in $\ell$, each left-child
step removes a right piece of length $R\ge\tau/|\ell|$.
After $n_L$ such steps $g_R\ge n_L\tau/|\ell|$.
At a further left-child step, \eqref{branching:phase-gap} gives
$n_LL<\kappa|\ell|$, and its split child has
$L>|\ell|/(\kappa+1)$. Hence there are at most
$\kappa(\kappa+1)$ left steps, and the same number of right steps.
A split-node chain contains at most $2\kappa(\kappa+1)+1$ nodes.
For any finite portion of the split-node forest inside $\ell$, terminal
split boxes are disjoint in their interiors and all longer than
$|\ell|/(\kappa+1)$, so there are at most $\kappa$ of them.
The forest is covered by their ancestor chains. There are at most
$c_\kappa$ split nodes inside $\ell$. This applies to every finite
partial tree, so also excludes an infinite forest.

Let $T_b$ be the ordinary minimizer tree stopped when $F\ge-b$.
Any split node of $S$ with $F<-b$ belongs to $T_b$, because all
ancestors also have $F<-b$ and the selections agree. There are at
most $4N-5$ such nodes by the one-dimensional theorem. Every remaining
node lies inside the first ancestor with $F\ge-b$, a leaf of $T_b$.
That tree has at most $4N-4$ leaves when $N\ge2$; each contains at
most $c_\kappa$ additional split nodes. When $N=1$, the root itself
is budget-$b$ valid and the forest bound applies directly.
\end{proof}

\begin{proof}[Proof of \cref{branching:sharp-induction}]
After a sharp coordinate is cut at its kink, each half has relaxation
value exactly zero: it contains a zero-valued kink, and the one-sided
hypothesis gives a nonnegative minimum. Its selected gap is zero.
It can never be selected on an invalid node while any positive gap
remains; some positive gap must remain because all $m_i\ge0$.

We prove the leaf bound by induction on the number $r$ of unsplit
sharp coordinates, allowing already-split sharp coordinates as above.
For $r=0$, the one-dimensional theorem gives at most $4N_A(\epsilon)$
leaves, including the case $N_A=1$.
For $r\ge1$, freeze the unsplit sharp intervals until the first sharp
split on each branch. Their gaps $\omega_j>0$ are constants and their
values are $-\omega_j$. Put
\[
 \tau=\max_j\omega_j,\quad
 \beta=\epsilon-\sum_j\omega_j,\quad
 b=\beta+r\tau\ge\epsilon.
\]
Every arbitrary-coordinate split in this initial phase satisfies
$F(A')<-\beta$ and $\omega(A')\ge\tau$, so its phase is a subtree
of the phase tree in \cref{branching:phase-lemma} with $\kappa=r$.
Since $N_A(b)\le N_A(\epsilon)=N$, that lemma bounds its internal
nodes by $A_rN$, also when $N_A(b)=1$. Its frontier partitions $A$
into $H\le(A_r+1)N$ intervals $A_j$.

Refine an optimal budget-$\epsilon$ partition of $A$ by these $H-1$
frontier breakpoints. Hereditary validity means every refined piece
is still valid. Therefore
\begin{equation}
 \sum_jN_{A_j}(\epsilon)\le N+H-1\le(A_r+2)N.
 \label{branching:frontier-budget}
\end{equation}
A frontier node is either valid or splits one sharp coordinate at
its kink. In the latter case its two children have $r-1$ unsplit sharp
coordinates and the same arbitrary interval $A_j$. Apply induction
to them. Overcounting valid frontier nodes by $H$ gives a final leaf
bound
\[
 H+2C_{r-1}\sum_jN_{A_j}(\epsilon)
 \le[(A_r+1)+2C_{r-1}(A_r+2)]N=C_rN.
\]
The phase lemma and induction prove termination as well.

Finally slice any original rectangular certificate at every sharp
kink. On a box meeting that slice, the other coordinate values are
zero and their gap terms are nonnegative, so its arbitrary-coordinate
interval is valid at budget $\epsilon$. These intervals cover $A$.
An interval cover by hereditary-valid intervals gives a partition
with no more pieces: repeatedly select, among covering intervals
containing the current endpoint, one with the farthest right endpoint,
and take the corresponding subinterval. Finite covers guarantee
progress, and no covering interval can be selected twice.
Thus $N_A(\epsilon)\le N_{\rm rect}(\epsilon)$, proving the claimed
node comparison. Every phase and selection in this proof uses the
fixed coordinate-wise oracle; sharpness for a different tied minimizer
cannot be inferred from sharpness for the selected one.
\end{proof}
